\chapter{Хэрэгжүүлэлт}

Энэхүү бүлэгт өмнөх бүлгүүдэд тодорхойлсон системийн шаардлага болон
архитектурын зохиомжийн дагуу хэрэгжүүлсэн annotation болон reflection
механизмд суурилсан dependency injection framework-ийн дэлгэрэнгүй
хэрэгжүүлэлтийг тайлбарлана. Хэрэгжүүлэлт нь давхаргат архитектурын
зарчмыг баримтлан annotation давхарга, scanner давхарга, container
давхарга, exception давхарга гэсэн үндсэн хэсгүүдэд хуваагдана. Үүнээс
гадна A1 сайжруулалтын ажлын хүрээнд тестийн хамралтыг өсгөх,
fail-fast exception-уудын мэдээллийн чанарыг сайжруулах, туршилтын
бүтцийг \codeword{unit} болон \codeword{integration} түвшинд цэгцлэх
ажлууд хийгдсэн. Иймээс энэ бүлэгт үндсэн хэрэгжүүлэлтээс гадна уг
сайжруулалтууд системийн шаардлага болон зохиомжтой хэрхэн уялдсан
байдгийг хамтад нь тайлбарлана.

%-------------------------------------------------------------
\section{Хэрэгжүүлэлтийн орчин ба бүтэц}
%-------------------------------------------------------------

Хэрэгжүүлэлтийг системийн шаардлагын бүлэгт тодорхойлсон технологийн
орчны дагуу Java~SE~17, Apache Maven~3.9.x, JUnit~5 орчинд гүйцэтгэв.
Гуравдагч талын нэмэлт сан огт ашиглаагүй бөгөөд зөвхөн JDK стандарт
сангийн Reflection API болон annotation боловсруулах механизмыг ашигласан.

Төслийн package бүтэц нь архитектурын давхаргатай шууд нийцэж байна.

\begin{verbatim}
mn.edu.num/
|-- annotation/
|   |-- Component.java
|   |-- Autowired.java
|   |-- Qualifier.java
|   `-- Scope.java
|-- container/
|   |-- ApplicationContext.java
|   |-- BeanDefinition.java
|   |-- BeanRegistry.java
|   |-- DependencyInjector.java
|   `-- ScopeType.java
|-- scanner/
|   `-- ClassPathScanner.java
|-- exception/
|   |-- BeanCreationException.java
|   |-- NoSuchBeanException.java
|   `-- CircularDependencyException.java
`-- app/
    |-- broken/, circular/, circular3/
    |-- conflict/, qualifer/, scope/
    `-- throwing/
\end{verbatim}

Энд \codeword{annotation}, \codeword{scanner}, \codeword{container},
\codeword{exception} багцууд нь framework-ийн үндсэн хэрэгжүүлэлт бөгөөд
\codeword{app} багц нь жишээ хэрэглээ болон тестийн scenario-г бүрдүүлэх
зориулалттай. A1 ажлын үед энэхүү бүтэц өөрчлөгдөөгүй, харин тус
бүтцийн зан авирыг баталгаажуулах тестийн сан мэдэгдэхүйц өргөжсөн.

\section{A1 сайжруулалтын ажлын хэрэгжилтийн хүрээ}

Нэмэлт хэрэгжүүлэлтийн баримтад A1 ажлыг ``одоогийн хэрэгжүүлэлтийг
сайжруулж, тестийн хамралтыг өсгөх'' гэж тодорхойлсон. Энэ ажлын
хүрээнд framework-ийн үндсэн логикийг шинэчлэхээс илүүтэйгээр дараах
гурван чиглэлд төвлөрсөн.

\begin{enumerate}
    \item \textbf{Branch-oriented тестийн баталгаажуулалт} --
    \codeword{ClassPathScanner}, \codeword{ApplicationContext},
    \codeword{BeanRegistry}, \codeword{DependencyInjector} зэрэг цөм
    классуудын салбарласан урсгал бүрийг нэгж тестээр шалгах;
    \item \textbf{Fail-fast exception-ийн мэдээллийн чанарыг
    сайжруулах} -- bean нэр, хайсан төрөл, боломжит candidate,
    үүсгэгдэж буй төлөв, үндсэн \textit{cause} зэрэг мэдээллийг
    exception мэдэгдэлд тусгах;
    \item \textbf{Тестийн бүтцийг цэгцлэх} -- framework-ийн нэгж
    зан авир болон бүхэл урсгалын шалгалтыг \codeword{unit} болон
    \codeword{integration} түвшинд салгаж, дахин ажиллуулахад илүү
    ойлгомжтой тестийн сан бүрдүүлэх.
\end{enumerate}

Ингэснээр A1 ажил нь 2-р бүлгийн шаардлагын биелэлтийг батлах, 3-р
бүлгийн layered architecture-ийн зөв зааглалыг хамгаалах, 4-р бүлгийн
хэрэгжүүлэлтийг ``ажиллаж буй загвар''-аас ``тестээр баталгаажсан
хэрэгжүүлэлт'' түвшинд ахиулах үүрэг гүйцэтгэсэн.

%-------------------------------------------------------------
\section{Annotation давхаргын хэрэгжүүлэлт}
%-------------------------------------------------------------

Annotation давхарга нь framework-ийн хэрэглэгчтэй харилцах гадаад
интерфэйсийн үүрэг гүйцэтгэнэ. Шаардлагын бүлэгт тодорхойлсон дөрвөн
custom annotation-г \codeword{@interface} түлхүүр үгээр тодорхойлсон бөгөөд
бүгд \codeword{@Retention(RetentionPolicy.RUNTIME)} тохиргоотой байна.
Энэ нь annotation-уудыг программ ажиллах үед Reflection API ашиглан унших
боломжийг бүрдүүлдэг бөгөөд container-ийн бүх дотоод логик үүнд
тулгуурлан ажилладаг.

%- - - - - - - - - - - - - - - - - - - - - - - - - - - - - -
\subsection{\texttt{@Component} annotation}
%- - - - - - - - - - - - - - - - - - - - - - - - - - - - - -

\codeword{@Component} annotation нь тухайн классыг IoC container-д bean
хэлбэрээр бүртгэхийг заах үндсэн тэмдэглэгээ юм. \codeword{@Target}
нь \codeword{ElementType.TYPE} байх тул зөвхөн класс болон интерфэйс
дээр тавигдана. \codeword{value()} атрибут нь bean-ийн нэрийг тодорхойлох
боломж олгох бөгөөд хоосон үед \codeword{ClassPathScanner} классын нэрийг
жижиг үсгээр эхлүүлж автоматаар нэр үүсгэнэ.

% ---- КОД ОРУУЛАХ ХЭСЭГ ----
\begin{lstlisting}[language=Java,
    caption={Component annotation-ийн хэрэгжүүлэлт},
    label={lst:component-annotation}]
package mn.edu.num.annotation;

import java.lang.annotation.*;
    
@Retention(RetentionPolicy.RUNTIME)
@Target(ElementType.TYPE)
@Documented
public @interface Component {
    String value() default "";
}

\end{lstlisting}

%- - - - - - - - - - - - - - - - - - - - - - - - - - - - - -
\subsection{\texttt{@Autowired} annotation}
%- - - - - - - - - - - - - - - - - - - - - - - - - - - - - -

\codeword{@Autowired} annotation нь field болон constructor дээр тавигдаж
тухайн хамаарлыг container автоматаар inject хийхийг заана.
\codeword{required} атрибут нь injection заавал байх эсэхийг тодорхойлох
бөгөөд \codeword{true} утгатай үед bean олдохгүй бол
\codeword{NoSuchBeanException} шидэгдэнэ.

% ---- КОД ОРУУЛАХ ХЭСЭГ ----
\begin{lstlisting}[language=Java,
    caption={\texttt{@Autowired} annotation-ийн хэрэгжүүлэлт},
    label={lst:autowired-annotation}]
package mn.edu.num.annotation;

import java.lang.annotation.ElementType;
import java.lang.annotation.Retention;
import java.lang.annotation.RetentionPolicy;
import java.lang.annotation.Target;

@Target({ElementType.FIELD, ElementType.CONSTRUCTOR})
@Retention(RetentionPolicy.RUNTIME)
public @interface Autowired {
    boolean required() default true;
}

\end{lstlisting}

%- - - - - - - - - - - - - - - - - - - - - - - - - - - - - -
\subsection{\texttt{@Qualifier} ба \texttt{@Scope} annotation}
%- - - - - - - - - - - - - - - - - - - - - - - - - - - - - -

\texttt{@Qualifier} annotation нь ижил интерфэйсийг хэрэгжүүлсэн олон
bean бүртгэгдсэн үед injection-д ашиглах bean-г нэрээр заах зорилготой.
Энэ нь \texttt{DependencyInjector} дотор тохирох bean хайх үед
\textit{ambiguity} үүсэхээс сэргийлнэ. \texttt{@Scope} annotation нь
bean-ийн амьдралын мөчлөгийг тодорхойлох бөгөөд \texttt{singleton}
болон \texttt{prototype} гэсэн хоёр утгыг дэмжинэ.

% ---- КОД ОРУУЛАХ ХЭСЭГ ----
\begin{lstlisting}[language=Java,
    caption={\texttt{@Qualifier} annotation-ийн хэрэгжүүлэлт},
    label={lst:qualifier-annotation}]
package mn.edu.num.annotation;

import java.lang.annotation.ElementType;
import java.lang.annotation.Retention;
import java.lang.annotation.RetentionPolicy;
import java.lang.annotation.Target;

@Target({ElementType.PARAMETER,  ElementType.FIELD})
@Retention(RetentionPolicy.RUNTIME)
public @interface Qualifier {
    String value();
}

\end{lstlisting}

\begin{lstlisting}[language=Java,
    caption={\texttt{@Scope} annotation-ийн хэрэгжүүлэлт},
    label={lst:scope-annotation}]
package mn.edu.num.annotation;

import java.lang.annotation.ElementType;
import java.lang.annotation.Retention;
import java.lang.annotation.RetentionPolicy;
import java.lang.annotation.Target;

@Target({ElementType.TYPE})
@Retention(RetentionPolicy.RUNTIME)
public @interface Scope {
    String value() default "singleton";
}


\end{lstlisting}
%-------------------------------------------------------------
\section{Exception давхаргын хэрэгжүүлэлт}
%-------------------------------------------------------------

Exception давхарга нь framework-ийн \textit{fail-fast} зарчмыг
хэрэгжүүлэх гол хэрэгсэл юм. Гурван тусгай exception класс нь
\texttt{RuntimeException}-г өвлөн тодорхойлогдсон тул хэрэглэгч
заавал барих шаардлагагүй боловч алдааны шалтгааныг тодорхой
мэдэгдлээр хүргэнэ.

\texttt{NoSuchBeanException} нь container-т хүссэн нэр эсвэл төрлийн
bean бүртгэгдээгүй үед шидэгдэнэ. A1 сайжруулалтын явцад уг exception-ийн
мэдэгдэлд хайсан bean-ийн нэр, төрөл, бүртгэлтэй bean-үүдийн жагсаалт,
олон candidate тохирсон үед \texttt{@Qualifier} ашиглах шаардлага зэрэг
оношлох мэдээллийг хавсаргах байдлаар баяжуулсан. \texttt{BeanCreationException}
нь Reflection API ашиглан bean объект үүсгэх явцад техникийн алдаа гарсан
үед \textit{cause} exception-г wrap хийж шидэгдэнэ. \texttt{CircularDependencyException}
нь A~$\rightarrow$~B~$\rightarrow$~A хэлбэрийн тойрог хамаарлыг container
эхлэх үед илрүүлэн шидэх бөгөөд үүсгэгдэж буй bean-үүдийн төлөвийг
мэдэгдэлд багтааснаар fail-fast шинж чанарыг илүү ойлгомжтой болгосон.

% ---- КОД ОРУУЛАХ ХЭСЭГ ----
\begin{lstlisting}[language=Java,
    caption={BeanCreationException классын хэрэгжүүлэлт},
    label={lst:beancreation-exception}]
package mn.edu.num.exception;

public class BeanCreationException extends RuntimeException{
    public BeanCreationException(String message, Throwable cause) {
        super(message, cause);
    }
}

\end{lstlisting}


\begin{lstlisting}[language=Java,
    caption={NoSuchBeanException классын хэрэгжүүлэлт},
    label={lst:nosuchbean-exception}]
package mn.edu.num.exception;

public class NoSuchBeanException extends RuntimeException {
    public NoSuchBeanException(String message) {
        super(message);
    }
}

\end{lstlisting}


\begin{lstlisting}[language=Java,
    caption={CircularDependencyException классын хэрэгжүүлэлт},
    label={lst:circulardependency-exception}]
package mn.edu.num.exception;

public class CircularDependencyException extends RuntimeException{
    public CircularDependencyException(String message){
        super(message);
    }
}


\end{lstlisting}
%-------------------------------------------------------------
\section{Container давхаргын хэрэгжүүлэлт}
%-------------------------------------------------------------

Container давхарга нь framework-ийн цөм бүрэлдэхүүн болох
\texttt{ScopeType} enum, \texttt{BeanDefinition}, \texttt{BeanRegistry}
болон \texttt{ApplicationContext} классуудаас бүрдэнэ.

%- - - - - - - - - - - - - - - - - - - - - - - - - - - - - -
\subsection{\texttt{ScopeType} ба \texttt{BeanDefinition}}
%- - - - - - - - - - - - - - - - - - - - - - - - - - - - - -

\texttt{ScopeType} нь bean-ийн амьдралын мөчлөгийг \texttt{String}
биш type-safe \texttt{enum} хэлбэрээр хадгалдаг. Энэ нь буруу утга
оруулах боломжийг compile цагт таслан зогсооно.

\texttt{BeanDefinition} нь bean-ийн бодит объект биш, харин тухайн
bean-г хэрхэн үүсгэх, ямар scope-той болохыг илэрхийлэх мета
мэдээллийн савлагаа юм. Дотор нь \texttt{beanName},
\texttt{beanClass}, \texttt{scope} гэсэн гурван талбар агуулна.

% ---- КОД ОРУУЛАХ ХЭСЭГ ----
\begin{lstlisting}[language=Java,
    caption={\texttt{ScopeType} enum классын хэрэгжүүлэлт},
    label={lst:scopetype-enum}]
package mn.edu.num.container;

public enum ScopeType {
    SINGLETON, PROTOTYPE;
}
\end{lstlisting}

\begin{lstlisting}[language=Java,
    caption={\texttt{BeanDefinition} классын хэрэгжүүлэлт},
    label={lst:beandefinition-class}]
package mn.edu.num.container;

public class BeanDefinition {
    private String name;
    private Class<?> beanClass;
    private ScopeType scope;

    public BeanDefinition(String name, Class<?> beanScope, ScopeType scope){
        this.name = name;
        this.beanClass = beanScope;
        this.scope = scope;
    }

    public String getBeanName() {return this.name;}
    public Class<?> getBeanClass() {return this.beanClass;}
    public ScopeType getScope() {return this.scope;}

    @Override
    public String toString() {
        return "BeanDefinition{" +
                "name='"  + name  + '\'' +
                ", class=" + beanClass.getSimpleName() +
                ", scope=" + scope +
                '}';
    }
}

\end{lstlisting}

%- - - - - - - - - - - - - - - - - - - - - - - - - - - - - -
\subsection{\texttt{BeanRegistry}}
%- - - - - - - - - - - - - - - - - - - - - - - - - - - - - -

\texttt{BeanRegistry} нь container-ийн санах ойн үүрэг гүйцэтгэнэ.
Дотроо хоёр \texttt{HashMap} агуулна: нэгдүгээрт,
bean-ийн нэрийг \texttt{BeanDefinition}-д харгалзуулан хадгалах
\texttt{definitions} Map, хоёрдугаарт, singleton bean-ийн бодит
instance-г cache хийх \texttt{singletonCache} Map. \texttt{findByType()}
метод нь \texttt{isAssignableFrom()} ашиглан интерфэйс болон
superclass-аар хайх боломжийг олгодог.

% ---- КОД ОРУУЛАХ ХЭСЭГ ----
\begin{lstlisting}[language=Java,
    caption={\texttt{BeanRegistry} классын хэрэгжүүлэлт},
    label={lst:beanregistry}]
package mn.edu.num.container;

import java.util.*;

public class BeanRegistry {
    // bean нэр → BeanDefinition
    private final Map<String, BeanDefinition> definitions = new HashMap<>();
    // singleton cache
    private final Map<String, Object> singletonCache = new HashMap<>();

    // Төрлөөр хайх
    public List<BeanDefinition> findByType(Class<?> type) {
        List<BeanDefinition> result = new ArrayList<>();
        for (BeanDefinition def : definitions.values()) {
            if (type.isAssignableFrom(def.getBeanClass())) {
                result.add(def);
            }
        }
        return result;
    }
}

\end{lstlisting}

%- - - - - - - - - - - - - - - - - - - - - - - - - - - - - -
\subsection{\texttt{ApplicationContext}}
%- - - - - - - - - - - - - - - - - - - - - - - - - - - - - -

\texttt{ApplicationContext} нь framework-ийн хамгийн чухал, төв
бүрэлдэхүүн хэсэг юм. Хэрэглэгч зөвхөн энэ классыг ашиглан
container-тай харилцана. Constructor-д \texttt{basePackage} нэр
дамжуулахад дараах дарааллаар startup явагдана.

\begin{enumerate}
    \item \texttt{ClassPathScanner}-ээр заасан package болон
          дэд package-уудыг нэгжиж \texttt{@Component} annotation-той
          классуудыг илрүүлнэ.
    \item Илэрсэн класс бүрийг \texttt{BeanDefinition} болгон
          \texttt{BeanRegistry}-д бүртгэнэ.
    \item \texttt{SINGLETON} scope-той bean бүрийг урьдчилан
          үүсгэж singleton cache - д хадгална.
    \item \texttt{DependencyInjector}-ээр \texttt{@Autowired}
          field бүрд тохирох dependency-г inject хийнэ.
\end{enumerate}

Тойрог хамааралыг илрүүлэхийн тулд \texttt{inCreation} гэсэн
\texttt{HashSet} ашиглана. Bean үүсгэж эхлэхэд нэрийг
\texttt{Set}-д нэмж, бэлэн болсны дараа хасна. Дахин тэр bean-г
үүсгэхийг оролдвол \texttt{Set}-д байгааг олж
\texttt{CircularDependencyException} шидэнэ.

A1 ажлын үед энэ хэсгийн чанарын баталгаажуулалт мэдэгдэхүйц
өргөжсөн. Тухайлбал singleton cache-д байгаа болон байхгүй
тохиолдол, prototype bean хүсэх үеийн урсгал, нэрээр болон төрлөөр
bean авах үеийн алдааны нөхцөл, олон candidate тохирох ambiguity,
мөн хоёр болон гурван шаттай circular dependency зэрэг нөхцөл бүрийг
тусгай тестээр баталгаажуулсан. Иймээс \texttt{ApplicationContext}
нь зөвхөн bean үүсгэх төв класс бус, мөн framework-ийн fail-fast
зан авирыг хэрэгжүүлэх үндсэн цэг болж байна.

% ---- КОД ОРУУЛАХ ХЭСЭГ ----
\begin{lstlisting}[language=Java,
    caption={\texttt{ApplicationContext} классын хэрэгжүүлэлт},
    label={lst:applicationcontext}]
package mn.edu.num.container;

import mn.edu.num.exception.BeanCreationException;
import mn.edu.num.exception.CircularDependencyException;
import mn.edu.num.exception.NoSuchBeanException;
import mn.edu.num.scanner.ClassPathScanner;

import java.util.*;

public class ApplicationContext implements DependencyInjector.ApplicationContextRef {

    private final BeanRegistry registry = new BeanRegistry();
    private final DependencyInjector injector = new DependencyInjector(registry);
    // Circular dependency илрүүлэх
    private final Set<String> inCreation = new HashSet<>();

    public ApplicationContext(String basePackage) {
        System.out.println("[Context] Эхэлж байна... package: " + basePackage);
        ClassPathScanner scanner = new ClassPathScanner(basePackage);
        List<BeanDefinition> definitions = scanner.scan();

        // Бүртгэх
        for (BeanDefinition def : definitions) {
            registry.register(def);
        }

        // Singleton bean-уудыг урьдчилан үүсгэх
        for (BeanDefinition def : definitions) {
            if (def.getScope() == ScopeType.SINGLETON) {
                getBean(def.getBeanName());
            }
        }
        System.out.println("[Context] Container бэлэн боллоо.");
    }

    @Override
    public Object getBean(String name) {
        BeanDefinition def = registry.getDefinition(name);

        // Кэш шалгахаас ӨМНӨ тойрог хамаарал шалгана
        if (inCreation.contains(name)) {
            throw new CircularDependencyException(
                    "Тойрог хамаарал илэрлээ: " + name
                            + " | Үүсгэгдэж байгаа: " + inCreation);
        }
        if (def.getScope() == ScopeType.SINGLETON) {
            if (registry.hasSingleton(name)) {
                return registry.getSingleton(name);
            }
            return createBean(def);
        } else {
            // Prototype — шинэ instance үүсгэнэ
            return createBean(def);
        }
    }
}
\end{lstlisting}

%-------------------------------------------------------------
\section{Scanner давхаргын хэрэгжүүлэлт}
%-------------------------------------------------------------

\texttt{ClassPathScanner} нь өгөгдсөн package болон түүний бүх
дэд package-уудыг рекурсив аргаар нэгжиж \texttt{@Component}
annotation-той классуудыг илрүүлэн \texttt{BeanDefinition} болгодог
модуль юм.

Ажиллах алгоритм нь дараах алхмуудаас бүрдэнэ. Эхлээд package - ийн нэрийн
цэг тэмдэгтийг зам тусгаарлагч болгон хөрвүүлж
(\texttt{"mn.edu.num"} $\rightarrow$ \texttt{"mn/edu/num"}),
\texttt{ClassLoader}-оос харгалзах \texttt{URL} авна. \texttt{URL}-г
\texttt{File} объект болгон рекурсив байдлаар нэгжиж, \texttt{.class}
өргөтгөлтэй файл бүрийг \texttt{Class.forName()} ашиглан JVM-д
ачаалдаг. Interface болон \texttt{\$} тэмдэгт агуулсан inner class-уудыг
шүүн гаргаж, \texttt{@Component} annotation байгаа классуудыг
\texttt{BeanDefinition} болгон буцаадаг.

A1 сайжруулалтын хүрээнд scanner давхаргын дараах edge case-уудыг
тусад нь шалгасан. Үүнд package олдохгүй байх, хоосон package scan хийх,
interface-ийг bean болгон бүртгэхгүй байх, custom bean name болон
decapitalize хийсэн default bean name үүсгэх, \texttt{@Scope("prototype")}
утгыг зөв унших, буруу scope утгыг singleton fallback болгон авч үзэх,
\texttt{.class} биш файлуудыг алгасах, \texttt{listFiles()} нь
\texttt{null} болох нөхцөл, мөн component биш классыг бүртгэхгүй байх
урьдал нөхцлүүд орно. Ингэснээр scanner-ийн алгоритм зөвхөн үндсэн
зам дээр бус, захын тохиолдлуудад ч тогтвортой ажиллаж байгааг баталсан.

% ---- КОД ОРУУЛАХ ХЭСЭГ ----
\begin{lstlisting}[language=Java,
    caption={\texttt{ClassPathScanner} классын хэрэгжүүлэлт},
    label={lst:classpathscanner}]
public class ClassPathScanner {

    private final String basePackage;

    public ClassPathScanner(String basePackage) {
        this.basePackage = basePackage;
    }
}
\end{lstlisting}

%-------------------------------------------------------------
\section{Injector давхаргын хэрэгжүүлэлт}
%-------------------------------------------------------------

\texttt{DependencyInjector} нь аль хэдийн үүсгэгдсэн bean объект
дээр \texttt{@Autowired} annotation-той field бүрд тохирох
dependency-г олж Reflection ашиглан оноодог модуль юм.

Inject хийх дараалал нь дараах байдалтай. Bean-ийн класс дахь
бүх field-ийг \texttt{getDeclaredFields()} ашиглан авч
\texttt{@Autowired} annotation байгааг шалгана. \texttt{@Qualifier}
annotation байвал заасан нэрээр \texttt{ApplicationContext}-оос
хайна. \texttt{@Qualifier} байхгүй бол \texttt{BeanRegistry.findByType()}
ашиглан төрлөөр хайж нэг candidate илэрсэн тохиолдолд inject хийнэ.
Олон candidate илэрвэл \textit{ambiguity} алдаа, нэг ч олдохгүй бол
\texttt{NoSuchBeanException} илгээнэ. Эцэст нь
\texttt{field.setAccessible(true)} ашиглан private field-д хандаж
\texttt{field.set(instance, dependency)} ашиглан утга оноодог.

\codeword{DependencyInjector.ApplicationContextRef} interface нь
\codeword{DependencyInjector} болон \codeword{ApplicationContext} хоёрын
хоорондох шууд хамаарлыг арилгах зорилготой. Энэ нь circular import
гарахаас сэргийлж, нэгж тест хийхэд mock объект ашиглах боломжийг
олгоно.

A1 ажлын үед injector давхаргын чанарын баталгаажуулалт дараах гол
урвалууд дээр төвлөрсөн. \texttt{@Qualifier} байхгүй үед dependency-г
төрлөөр нь resolve хийх, \texttt{@Qualifier} өгөгдсөн үед яг заасан
bean-ийг нэрээр нь сонгох, dependency байхгүй үед тайлбарлагдсан
\texttt{NoSuchBeanException} шидэх, олон candidate илэрсэн үед ambiguity
алдаа өгөх, мөн reflection-оор field утга оноох явцад гарч болох
\texttt{IllegalAccessException}-ийг \texttt{BeanCreationException}-д wrap
хийх зэрэг урсгалуудыг нэгж тестээр баталгаажуулсан. Энэ нь injector
давхаргыг бодит нөхцөлд илүү найдвартай болгосон.

% ---- КОД ОРУУЛАХ ХЭСЭГ ----
\begin{lstlisting}[language=Java,
    caption={\texttt{DependencyInjector} классын хэрэгжүүлэлт},
    label={lst:dependencyinjector}]
public class DependencyInjector {

    private final BeanRegistry registry;

    public DependencyInjector(BeanRegistry registry) {
        this.registry = registry;
    }
}
\end{lstlisting}

%-------------------------------------------------------------
\section{Туршилт}
%-------------------------------------------------------------

Хэрэгжүүлсэн framework-ийн үндсэн функциональ шаардлага бүрийг
JUnit~5 тестийн орчинд баталгаажуулав. A1 ажлын үед тестийн
стратегийг \codeword{unit} болон \codeword{integration} гэсэн хоёр түвшинд
салгаж, framework-ийн цөм классуудын салбарласан урсгал болон
бүхэл startup--lookup--inject урсгалыг тус тусад нь шалгах зохион
байгуулалт бүрдүүлсэн.

\subsection{Нэгж тестийн бүтэц ба шалгасан нөхцөлүүд}

Нэгж тестүүд нь тухайн нэг класс эсвэл туслах логикийн зан авирыг
тусгаарлан шалгахад чиглэсэн. Энэ түвшинд \texttt{BeanRegistryTest},
\texttt{ClassPathScannerTest}, \texttt{DependencyInjectorTest},
\texttt{ExceptionTest} гэсэн дөрвөн тест класс ашиглагдсан.

\begin{itemize}
    \item \texttt{BeanRegistryTest} -- definition бүртгэх, singleton cache,
    төрлөөр хайх, давхардсан bean нэрийг overwrite хийх, \texttt{toString()}
    үр дүнг шалгах;
    \item \texttt{ClassPathScannerTest} -- package байхгүй, хоосон package,
    interface ignore хийх, custom/default bean name, prototype scope,
    component биш класс, \texttt{null} болон хоосон текстийн edge case-уудыг шалгах;
    \item \texttt{DependencyInjectorTest} -- dependency-г төрлөөр болон
    \texttt{@Qualifier}-аар resolve хийх, dependency байхгүй болон ambiguity
    үүсэх үед exception шидэх, reflection field set-ийн алдааг wrap хийх;
    \item \texttt{ExceptionTest} -- exception constructor-уудын мэдэгдэл,
    \texttt{ApplicationContext} болон \texttt{ClassPathScanner}-ийн хүрэхэд
    хэцүү branch-уудыг reflection ашиглан баталгаажуулах.
\end{itemize}

Эдгээр нэгж тестүүд нь A1 ажлын ``branch бүрийг илрүүлж тестлэх''
зорилтыг хэрэгжүүлж, framework-ийн цөм логикийн захын тохиолдлуудыг
системтэйгээр хамгаалж өгсөн.

\subsection{Интеграцийн тестийн бүтэц ба A1 ажлын үр дүн}

Интеграцийн тестүүд нь \texttt{ApplicationContext}-г бүхэлд нь
ажиллуулж, scan, register, inject, getBean гэсэн дамжлагууд бодит
жишээ package-ууд дээр хэрхэн ажиллаж байгааг шалгана. Энэ түвшинд
\texttt{BeanCreationIntegrationTest}, \texttt{BeanLookupIntegrationTest},
\texttt{DependencyInjectionIntegrationTest},
\texttt{QualifierIntegrationTest}, \texttt{ScopeIntegrationTest},
\texttt{CircularDependencyIntegrationTest} гэсэн зургаан тест класс
ашиглагдсан.

\clearpage
\renewcommand{\arraystretch}{1.25}
\footnotesize
\setlength{\tabcolsep}{4pt}
\begin{table}[p]
\centering
\caption{A1 ажлын хүрээнд ашигласан тестийн матриц}
\label{tab:a1-test-matrix}
\begin{tabular}{|L{0.23\textwidth}|L{0.13\textwidth}|L{0.10\textwidth}|L{0.50\textwidth}|}
\hline
\textbf{Тест класс} & \textbf{Түвшин} & \textbf{Тестийн тоо} & \textbf{Шалгасан үндсэн зан авир} \\ \hline
\shortstack[l]{\codeword{BeanRegistry}\\\codeword{Test}} & Unit & 8 & \shortstack[l]{Definition бүртгэл, cache,\\төрлөөр хайлт, overwrite,\\\codeword{toString()}} \\ \hline
\shortstack[l]{\codeword{ClassPathScanner}\\\codeword{Test}} & Unit & 12 & \shortstack[l]{Package scan, custom/default нэр,\\scope, interface ignore, edge case} \\ \hline
\shortstack[l]{\codeword{DependencyInjector}\\\codeword{Test}} & Unit & 6 & \shortstack[l]{Inject by type, qualifier, missing\\dependency, ambiguity, field access error} \\ \hline
\codeword{ExceptionTest} & Unit & 5 & \shortstack[l]{Exception constructor,\\reflection-оор өдөөсөн branch шалгалт} \\ \hline
\shortstack[l]{\codeword{BeanCreation}\\\codeword{IntegrationTest}} & Integration & 3 & \shortstack[l]{Default constructor байхгүй,\\constructor exception} \\ \hline
\shortstack[l]{\codeword{BeanLookup}\\\codeword{IntegrationTest}} & Integration & 3 & \shortstack[l]{Нэрээр болон төрлөөр lookup fail,\\multiple beans ambiguity} \\ \hline
\shortstack[l]{\codeword{DependencyInjection}\\\codeword{IntegrationTest}} & Integration & 2 & \shortstack[l]{\codeword{EmailService}, \codeword{OrderService}\\dependency inject} \\ \hline
\shortstack[l]{\codeword{Qualifier}\\\codeword{IntegrationTest}} & Integration & 3 & \shortstack[l]{\codeword{@Qualifier} resolve,\\custom bean name lookup} \\ \hline
\shortstack[l]{\codeword{Scope}\\\codeword{IntegrationTest}} & Integration & 3 & \shortstack[l]{Singleton, prototype,\\invalid scope fallback} \\ \hline
\shortstack[l]{\codeword{CircularDependency}\\\codeword{IntegrationTest}} & Integration & 2 & \shortstack[l]{Хоёр болон гурван шаттай\\circular dependency илрүүлэлт} \\ \hline
\textbf{Нийт} & -- & \textbf{47} & Framework-ийн гол болон захын урсгалуудыг хамарсан \\ \hline
\end{tabular}
\end{table}
\normalsize
\renewcommand{\arraystretch}{1}

Хүснэгт~\ref{tab:a1-test-matrix}-аас харахад A1 ажлын хүрээнд тестийн
сан зөвхөн тоогоор өссөнгүй, харин framework-ийн давхарга тус бүрийн
хариуцлагад уялдсан бүтэцтэй болсон. Ингэснээр regression эрсдэлийг
бууруулах, өөрчлөлтийн нөлөөг хурдан илрүүлэх, дараагийн A2--A4 ажлуудын
суурь чанарын хамгаалалтыг бүрдүүлэх боломжтой болсон.

\subsection{\texttt{mvn test} гүйцэтгэлийн баталгаа}

Төслийн үндсэн директорт \codeword{mvn test} командыг ажиллуулж шалгахад
framework-ийн тестийн сан амжилттай гүйцэтгэсэн. Үүний үр дүнд нийт 47
тест ажиллаж, \codeword{0 failures}, \codeword{0 errors}, \codeword{0 skipped}
байдлаар build success болсон. Энэ нь A1 ажлын тайлбар зөвхөн кодын
түвшний дүгнэлт бус, бодит build ба тестийн гүйцэтгэлийн үр дүнд
тулгуурлаж байгааг харуулна.

%-------------------------------------------------------------
\section{Бүлгийн дүгнэлт}
%-------------------------------------------------------------

Энэ бүлэгт системийн шаардлага болон архитектурын зохиомжийн дагуу
annotation болон reflection механизмд суурилсан dependency injection
framework-ийн туршилтын загвар болгон хэрэгжүүллээ. Давхаргат архитектурын зарчмыг
баримтлан annotation, scanner, container, injector, exception гэсэн
таван давхаргыг тус тусад нь хэрэгжүүлж, тэдгээрийн хоорондын
хамаарлыг интерфэйсээр зохицуулав.

Хэрэгжүүлэлтийн гол онцлог нь зөвхөн Java~SE стандарт сангийн
Reflection API болон annotation inspection механизмыг ашигласнаар
framework-ийн дотоод ажиллагааг гуравдагч талын сангаас бүрэн
хамааралгүй байдлаар харуулсан явдал юм. Мөн A1 сайжруулалтын ажлын
үр дүнд fail-fast exception-уудын мэдээллийн чанар сайжирч,
\texttt{unit} болон \texttt{integration} түвшний нийт 10 тест класс,
47 тестийн тохиолдлоор framework-ийн гол болон захын урсгалуудыг
баталгаажууллаа. Ингэснээр энэхүү хэрэгжүүлэлт нь зөвхөн ажилладаг
туршилтын загвар бус, харин архитектур, шаардлага, туршилтын нотолгоо
хоорондоо уялдсан чанарын баталгаатай хэрэгжүүлэлт болж байна.
